\documentclass[11pt]{article}

\usepackage[margin=1in]{geometry}
\usepackage{amsmath,amssymb,amsthm}
\usepackage{enumitem}
\usepackage{hyperref}
\usepackage{xcolor}
\usepackage{booktabs}

\hypersetup{colorlinks=true, linkcolor=blue, urlcolor=blue, citecolor=blue}

\title{TOAE209/TOAH209 -- Design and Analysis of Algorithms\\[4pt]
\large Week 11: Theoretical Questions}
\author{Foreign Trade University}
\date{}

\begin{document}
\maketitle

\section*{Instructions}

Part A contains 17 questions requiring written explanations or short proofs
(including the key proofs of Hall's and K\H{o}nig's theorems). Part B is a
10-question quiz (true/false and multiple choice). Attempt every question on
your own before consulting Part C (Solutions) at the end of this document.

\section*{Part A: Theory Questions}

\begin{enumerate}[label=\textbf{A\arabic*.}, leftmargin=2.2em]

\item Define a \emph{bipartite graph} in two equivalent ways: via a vertex
2-coloring, and via cycle lengths. Briefly explain why a BFS/DFS 2-coloring
algorithm (Problem 1) detects non-bipartiteness exactly when it finds an edge
between two same-colored vertices.

\item State the Maximum Bipartite Matching problem precisely (input, what a
matching is, objective). Then describe the standard reduction to maximum flow:
which vertices and edges are added, and what capacities they receive.

\item Prove that, in the matching-to-flow reduction, the maximum matching size
equals the maximum $s$--$t$ flow value. Where is the \emph{integrality theorem}
used, and why is it essential?

\item Kuhn's algorithm (Problem 2) repeatedly searches for an \emph{augmenting
path}. Define an augmenting path with respect to a matching $M$, and explain why
finding one lets you increase the matching size by exactly one. Why does the
algorithm terminate, and how many augmentations occur in the worst case?

\item Define the \emph{neighborhood} $N(S)$ of a set $S \subseteq X$ and a
\emph{deficient} set. State \textbf{Hall's theorem} precisely.

\item Prove the ``easy'' (necessity) direction of Hall's theorem: if a matching
saturates all of $X$, then $|N(S)| \ge |S|$ for every $S \subseteq X$.

\item Prove the ``hard'' (sufficiency) direction of Hall's theorem using
max-flow min-cut: if every $S \subseteq X$ satisfies $|N(S)| \ge |S|$, then an
$X$-saturating matching exists. (Hint: show the minimum cut of the matching
network equals $\min_{S}\,(|X| - |S| + |N(S)|)$.)

\item When no perfect matching of $X$ exists, Problem 4 returns a deficient set
$S$ as a \emph{certificate}. Explain why a deficient set is a convincing,
easily-checkable proof that no $X$-saturating matching can exist.

\item Define a \emph{vertex cover}. Prove the weak-duality bound: in \emph{any}
graph, $|M| \le |C|$ for every matching $M$ and vertex cover $C$.

\item State \textbf{K\H{o}nig's theorem}. Give a concrete example of a
\emph{non-bipartite} graph where maximum matching $<$ minimum vertex cover
(so K\H{o}nig fails), and report both values.

\item Prove K\H{o}nig's theorem constructively: starting from a maximum matching
$M$, build a vertex cover of size $|M|$ using alternating paths from the
unmatched $X$-vertices. (This is exactly the construction implemented in
Problem 5.)

\item Explain how Hall's theorem can be derived as a corollary of K\H{o}nig's
theorem.

\item State \textbf{Menger's theorem} (edge version) and explain the reduction
(Problem 7) that computes the maximum number of edge-disjoint $s$--$t$ paths.
Why does using capacity-$1$ edges, together with the integrality theorem,
guarantee the paths are edge-disjoint?

\item Describe the \emph{node-splitting} gadget (Problem 8) for computing
\emph{vertex-disjoint} paths. Give a small example of a graph in which the number
of edge-disjoint $s$--$t$ paths is strictly greater than the number of
vertex-disjoint $s$--$t$ paths.

\item State the \emph{circulation with demands and lower bounds} problem. Explain
the two-step reduction to max flow (Problem 11): how lower bounds are eliminated,
and how a super-source/super-sink decide feasibility. What is the role of the
condition $\sum_v d_v = 0$?

\item Describe the \emph{project selection / maximum-weight closure} problem and
its min-cut reduction (Problem 10). Why are the prerequisite edges given
\emph{infinite} capacity, and what does the source side of the minimum cut
represent?

\item Both \textbf{baseball elimination} (Problem 9) and \textbf{image
segmentation} are solved with flow/cut. For each, state in one or two sentences
what the flow/cut value represents and how the answer is read off. What do these
two applications, together with project selection, have in common at the level of
the reduction technique?

\end{enumerate}

\newpage
\section*{Part B: Quiz}

\begin{enumerate}[label=\textbf{B\arabic*.}, leftmargin=2.2em]

\item \textbf{(True/False)} A graph is bipartite if and only if it contains no
odd-length cycle.

\item \textbf{(Multiple choice)} In the standard reduction of bipartite matching
to max flow, the edges between the two sides $X$ and $Y$ are given capacity:
\begin{enumerate}[label=(\alph*)]
    \item $0$
    \item $1$
    \item $\infty$
    \item the degree of the endpoint in $X$
\end{enumerate}

\item \textbf{(Short answer)} A bipartite graph has $X = \{x_1, x_2\}$,
$Y = \{y_1\}$, with edges $x_1 y_1$ and $x_2 y_1$. Does $X$ have a perfect
matching? If not, give a deficient set $S$ and the values $|S|$ and $|N(S)|$.

\item \textbf{(True/False)} In every graph, the size of a maximum matching equals
the size of a minimum vertex cover.

\item \textbf{(Multiple choice)} K\H{o}nig's theorem states that maximum matching
equals minimum vertex cover in:
\begin{enumerate}[label=(\alph*)]
    \item all graphs
    \item bipartite graphs
    \item planar graphs
    \item trees only
\end{enumerate}

\item \textbf{(True/False)} To count \emph{vertex-disjoint} $s$--$t$ paths, one
splits each internal vertex $v$ into $v_{\text{in}} \to v_{\text{out}}$ with
capacity $1$.

\item \textbf{(Short answer)} In the project-selection min-cut network, why must
the prerequisite edges have infinite capacity?

\item \textbf{(Multiple choice)} In baseball elimination, team $x$ is declared
\emph{not eliminated} when the maximum flow:
\begin{enumerate}[label=(\alph*)]
    \item is zero
    \item saturates every edge \emph{into} the sink $t$
    \item saturates every game edge \emph{out of} the source $s$
    \item equals the number of teams
\end{enumerate}

\item \textbf{(Short answer)} In a circulation with demands, after eliminating an
edge's lower bound $\ell$ on edge $(u,v)$, how do the demands of $u$ and $v$
change, and how does the edge's capacity change?

\item \textbf{(True/False)} Image segmentation (foreground/background labeling
with separation penalties) can be solved by computing a minimum $s$--$t$ cut.

\end{enumerate}

\newpage
\section*{Part C: Solutions}

\subsection*{Part A}

\begin{enumerate}[label=\textbf{A\arabic*.}, leftmargin=2.2em]

\item A graph is \textbf{bipartite} if its vertices split into $X \cup Y$ with
every edge crossing between $X$ and $Y$; equivalently, if the vertices can be
2-colored (say colors $0$ and $1$) so that adjacent vertices get different
colors. A second equivalent characterization: a graph is bipartite iff it has
\emph{no odd-length cycle}. The BFS 2-coloring algorithm colors a start vertex
$0$ and alternates colors across edges; if it ever finds an edge whose endpoints
already have the \emph{same} color, the two BFS paths to those endpoints, plus
that edge, close an \emph{odd} cycle -- so the graph cannot be bipartite. If no
such edge is ever found, the coloring is a valid bipartition.

\item \textbf{Input:} a bipartite graph $G = (X \cup Y, E)$. A \textbf{matching}
$M \subseteq E$ is a set of edges no two of which share an endpoint. The
\textbf{objective} is to find a matching of maximum size. \textbf{Reduction to
flow:} add a source $s$ and sink $t$; add $s \to x$ of capacity $1$ for each
$x \in X$; add $y \to t$ of capacity $1$ for each $y \in Y$; orient each original
edge as $x \to y$ with capacity $1$. The maximum $s$--$t$ flow value equals the
maximum matching size.

\item \emph{Matching $\Rightarrow$ flow:} a matching $M$ of size $k$ yields a
flow of value $k$ by routing one unit along $s \to x \to y \to t$ for each
$(x,y) \in M$; disjointness of $M$ keeps every capacity-$1$ edge within its
capacity. \emph{Flow $\Rightarrow$ matching:} by the \textbf{integrality
theorem}, since all capacities are integers there is a maximum flow that is
\emph{integral}, so each edge carries $0$ or $1$ unit. The saturated middle edges
form a set $M$; each $x$ receives at most $1$ unit from $s$, so it sends at most
$1$ unit onward, i.e.\ $x$ is in at most one edge of $M$ (likewise each $y$). So
$M$ is a matching with $|M| = $ flow value. Integrality is \emph{essential}:
without it a max flow could split a unit fractionally across two middle edges,
which would not correspond to any matching.

\item An \textbf{augmenting path} (w.r.t.\ $M$) is a path that starts and ends at
\emph{free} vertices and \emph{alternates} between edges not in $M$ and edges in
$M$. Because the two endpoints are free and the path alternates, it contains one
more non-matching edge than matching edge; \emph{flipping} the path (removing its
matching edges from $M$ and adding its non-matching edges) yields a valid
matching with exactly one more edge. Kuhn's algorithm tries to find an augmenting
path from each $X$-vertex; each success increases $|M|$ by $1$, and since
$|M| \le \min(|X|,|Y|)$ there are at most that many augmentations, after which no
augmenting path remains and $M$ is maximum (Berge's theorem).

\item For $S \subseteq X$, $N(S) = \{y \in Y : y \text{ adjacent to some }
x \in S\}$. $S$ is \textbf{deficient} if $|N(S)| < |S|$. \textbf{Hall's theorem:}
a bipartite graph has a matching saturating all of $X$ if and only if
$|N(S)| \ge |S|$ for every $S \subseteq X$ (i.e.\ no deficient set exists).

\item Suppose a matching $M$ saturates every vertex of $X$. Fix any
$S \subseteq X$. Each $x \in S$ is matched to a distinct partner $\mu(x) \in Y$,
and $\mu(x) \in N(S)$ since $x \in S$ is adjacent to it. The map $x \mapsto
\mu(x)$ is injective (a matching pairs distinct $x$'s with distinct $y$'s), so
$N(S)$ contains at least $|S|$ distinct vertices: $|N(S)| \ge |S|$.

\item Build the matching network. An $X$-saturating matching exists iff the
maximum flow value is $|X|$, iff (max-flow min-cut) the minimum cut capacity is
$|X|$. Take any finite-capacity $s$--$t$ cut $(A,B)$, $s\in A$. Let
$S = X \cap B$. A finite cut cannot include any infinite-capacity edge; here all
edges are finite, but the \emph{cheapest} cut that ``abandons'' $S$ to the sink
side must (i) cut the edges $s \to x$ for each $x \in X \cap A = X \setminus S$,
costing $|X| - |S|$, and (ii) cut the edges $y \to t$ for every $y \in N(S)$ that
it keeps on the source side, costing $|N(S)|$ (it must keep $N(S)$ on the source
side, because each middle edge $x \to y$ with $x \in S$ would otherwise be a
source-to-sink crossing it has to pay for, and putting $N(S)$ on the source side
and cutting $y \to t$ is at least as cheap). Hence the minimum cut equals
$\min_{S \subseteq X}\big(|X| - |S| + |N(S)|\big)$. This is $< |X|$ exactly when
some $S$ has $|N(S)| < |S|$. So if \emph{no} set is deficient, the minimum cut is
$|X|$, the max flow is $|X|$, and an $X$-saturating matching exists.

\item A deficient set $S$ has $|N(S)| < |S|$. Any matching can match the vertices
of $S$ \emph{only} to vertices of $N(S)$ (those are their only neighbors), and a
matching pairs them injectively; but there are fewer than $|S|$ available
partners in $N(S)$, so at least one vertex of $S$ must go unmatched. Hence no
matching can saturate $X$. The certificate is \emph{easy to check}: given $S$,
one computes $N(S)$ in $O(E)$ time and verifies $|N(S)| < |S|$ -- no need to
re-run any matching algorithm.

\item A \textbf{vertex cover} $C$ is a set of vertices such that every edge has
at least one endpoint in $C$. \emph{Weak duality:} let $M$ be any matching and
$C$ any vertex cover. Each edge of $M$ must be covered, so it has an endpoint in
$C$. Distinct edges of $M$ share no endpoint, so they require \emph{distinct}
cover vertices. Therefore $|C| \ge |M|$. Taking the maximum matching and minimum
cover gives $\max$-matching $\le$ $\min$-cover.

\item \textbf{K\H{o}nig's theorem:} in a \emph{bipartite} graph, maximum matching
size $=$ minimum vertex cover size. \emph{Non-bipartite counterexample:} the
triangle $K_3$ on vertices $\{1,2,3\}$ with edges $\{12,23,31\}$. Its maximum
matching has size $1$ (any two edges share a vertex), but its minimum vertex
cover has size $2$ (one vertex covers only two of the three edges). So
$1 < 2$, and K\H{o}nig fails for this non-bipartite graph. (This is precisely the
example checked in Problem 12.)

\item Let $M$ be a maximum matching. Let $U \subseteq X$ be the unmatched
$X$-vertices. \emph{Mark} every vertex reachable from $U$ by an alternating path
(non-matching edge $X\to Y$, then matching edge $Y \to X$, repeating); call the
marked set $Z$. Define $C = (X \setminus Z) \cup (Y \cap Z)$.
\emph{$C$ covers every edge:} if some edge $(x,y)$ had $x \in X \cap Z$ and
$y \in Y \setminus Z$, then either $(x,y) \notin M$ (and we could extend the
alternating path from the marked $x$ to mark $y$ -- contradiction) or $(x,y) \in M$
(and $x$ could only have been marked via its matching partner $y$, so $y$ is
marked -- contradiction). \emph{$|C| = |M|$:} every vertex of $C$ is matched (an
unmatched $X$-vertex lies in $U \subseteq Z$ so is excluded from $X \setminus Z$;
an unmatched $Y$-vertex in $Z$ would give an augmenting path, contradicting
maximality), and no matching edge has both endpoints in $C$. Hence each matching
edge contributes exactly one endpoint to $C$, giving a bijection $M
\leftrightarrow C$ and $|C| = |M|$. Combined with weak duality ($|M| \le |C|$),
$C$ is a minimum cover of size exactly the maximum matching.

\item By K\H{o}nig's theorem, the minimum vertex cover of a bipartite graph has
size equal to the maximum matching. $X$ has a perfect matching iff the maximum
matching size is $|X|$, iff the minimum vertex cover has size $|X|$. One then
shows directly that a vertex cover smaller than $|X|$ exists iff some $S
\subseteq X$ is deficient: a cover that omits some $x \in S$ must include all of
$N(\{x\})$, and aggregating over a deficient $S$ produces a cover of size
$< |X|$; conversely a small cover yields a deficient set. Thus ``no perfect
matching'' $\iff$ ``a deficient set exists'' -- which is Hall's theorem.

\item \textbf{Menger's theorem (edge version):} the maximum number of
edge-disjoint $s$--$t$ paths equals the minimum number of edges whose removal
disconnects $t$ from $s$. \emph{Reduction:} give every edge capacity $1$ and
compute a max flow. By the integrality theorem the max flow is integral, so each
edge carries $0$ or $1$ unit; the positive-flow edges decompose into $s$--$t$
paths, and since each edge carries at most one unit, \emph{no edge is shared} --
the paths are edge-disjoint. The number of such paths equals the flow value,
which by max-flow min-cut equals the minimum edge cut.

\item \textbf{Node-splitting:} replace each internal vertex $v$ by two vertices
$v_{\text{in}}, v_{\text{out}}$ joined by an edge of capacity $1$; reroute every
original edge $(u,v)$ as $u_{\text{out}} \to v_{\text{in}}$ (capacity $1$). The
capacity-$1$ internal edge forces at most one path through $v$, so a max flow
counts \emph{vertex}-disjoint paths. \emph{Example:} vertices $s, a, b, m, t$
with edges $s\to a, s\to b, a\to m, b\to m, m\to t$. There are $2$
\emph{edge}-disjoint paths ($s\,a\,m\,t$ and $s\,b\,m\,t$ -- they share no edge),
but only $1$ \emph{vertex}-disjoint path, since both must pass through the single
internal vertex $m$. (This is the contrast tested between Problems 7 and 8.)

\item \textbf{Problem:} each vertex $v$ has a demand $d_v$ ($>0$ absorb, $<0$
supply); each edge has a lower bound $\ell_e$ and capacity $c_e$; we seek a flow
$f_e \in [\ell_e, c_e]$ with net flow into each $v$ equal to $d_v$.
\emph{Reduction:} (1) \emph{eliminate lower bounds}: for each edge $(u,v)$ with
bound $[\ell, c]$, set its capacity to $c - \ell$ and update $d_u \mathrel{+}=
\ell$, $d_v \mathrel{-}= \ell$ (we have committed to push $\ell$ units). (2) Add
a super-source $S$ and super-sink $T$: for each $v$ with adjusted $d_v < 0$, add
$S \to v$ of capacity $-d_v$; for each $v$ with adjusted $d_v > 0$, add $v \to T$
of capacity $d_v$. A feasible circulation exists iff the max $S$--$T$ flow
saturates all $S$-edges (equivalently equals the total positive demand). The
condition $\sum_v d_v = 0$ (after adjustment, total supply $=$ total demand) is
\emph{necessary}: flow conservation forces everything supplied to be absorbed,
so unequal totals make feasibility impossible -- Problem 11 checks this balance
before (or alongside) the flow computation.

\item \textbf{Project selection:} each project $v$ has profit $p_v$ (possibly
negative); a prerequisite $(p,q)$ requires that selecting $p$ forces selecting
$q$; we seek a prerequisite-closed set $S$ maximizing $\sum_{v\in S} p_v$.
\emph{Min-cut reduction:} $s \to v$ with capacity $p_v$ for $p_v > 0$;
$v \to t$ with capacity $-p_v$ for $p_v < 0$; prerequisite edges $p \to q$ with
\emph{infinite} capacity. The infinite capacities ensure no \emph{finite} cut
ever separates $p$ (source side) from a prerequisite $q$ (sink side) -- so the
source side is always a \emph{closed} set, i.e.\ feasible. The source side of the
minimum cut is the optimal project set, and the maximum profit is
$\big(\sum_{p_v>0} p_v\big) - (\text{min cut})$.

\item \textbf{Baseball elimination:} the max flow equals the number of remaining
games (among the other teams) that can be distributed without any team exceeding
$x$'s best-possible win total $W$. Team $x$ is \emph{not eliminated} iff the flow
\emph{saturates all game edges} (every remaining game can be ``played'' in the
model); otherwise it is eliminated, and the minimum cut names a set of teams that
collectively must accumulate too many wins. \textbf{Image segmentation:} the
minimum $s$--$t$ cut's capacity equals the total cost
$\sum_{i\in F} b_i + \sum_{j \in B} a_j + \sum_{\text{split}} p_{ij}$, and the
source/sink sides of the cut give the foreground/background labeling.
\emph{Common technique:} all three (baseball, project selection, segmentation)
are \textbf{min-cut} reductions -- they encode a combinatorial trade-off as the
capacity of a cut, then read the optimal partition off the source side. (Baseball
is most naturally posed as a \emph{saturation} check, which is the max-flow side
of the same coin.)

\end{enumerate}

\subsection*{Part B}

\begin{enumerate}[label=\textbf{B\arabic*.}, leftmargin=2.2em]

\item \textbf{True.} A graph is bipartite iff it can be 2-colored iff it has no
odd cycle -- these are three equivalent characterizations.

\item \textbf{(b) $1$.} The middle edges $x \to y$ get capacity $1$; the
capacity-$1$ edges from $s$ and into $t$ then enforce that each vertex is matched
at most once. (Using capacity $\infty$ on middle edges also works since the
endpoint edges already cap usage, but capacity $1$ is the standard choice.)

\item \textbf{No}, $X$ has no perfect matching. Both $x_1$ and $x_2$ have their
only neighbor $y_1$, so $S = \{x_1, x_2\}$ is deficient: $|S| = 2$ but
$|N(S)| = |\{y_1\}| = 1 < 2$.

\item \textbf{False.} This equality (K\H{o}nig's theorem) holds only for
bipartite graphs. In general, $\max$-matching $\le \min$-cover, and the triangle
$K_3$ shows the inequality can be strict ($1 < 2$).

\item \textbf{(b) bipartite graphs.} K\H{o}nig's theorem is specifically about
bipartite graphs; the analogous statement fails for general graphs.

\item \textbf{True.} The internal edge $v_{\text{in}} \to v_{\text{out}}$ of
capacity $1$ limits the flow through $v$ to one unit, forcing the paths to be
vertex-disjoint.

\item Infinite capacity guarantees that \emph{no finite-capacity cut} ever cuts a
prerequisite edge, so a project $p$ on the source side can never be separated
from its prerequisite $q$. This forces the source side of any minimum (finite)
cut to be \emph{closed} under prerequisites -- i.e.\ a feasible project set.

\item \textbf{(c) saturates every game edge out of the source $s$.} If all
remaining inter-team games can be routed to winners without pushing any team past
$W = w_x + (\text{remaining}_x)$, then $x$ can still finish (tied for) first.

\item The capacity becomes $c - \ell$ (we ``pre-pushed'' $\ell$ units), the
demand of $u$ increases by $\ell$ ($d_u \mathrel{+}= \ell$, since $u$ must now
also account for sending those $\ell$ units), and the demand of $v$ decreases by
$\ell$ ($d_v \mathrel{-}= \ell$, since $v$ already receives those $\ell$ units).

\item \textbf{True.} Foreground/background labeling with per-pixel likelihoods
and pairwise separation penalties is exactly a minimum $s$--$t$ cut: $s$ is the
foreground terminal, $t$ the background terminal, and the cut capacity equals the
total cost being minimized.

\end{enumerate}

\end{document}
